\section{The Creation Model}

Before anything, there was only the \textbf{source}. It had no edge and no end. It
was not a being you could point to. It filled everything there was, completely, so
that there was no empty place anywhere. There was nothing it could be set against,
because there was nothing else at all.

The source wished to give. Not because it lacked anything, but because it is the
nature of goodness to pour itself out. To give, there must be someone to receive.
So there had to be a world, and someone in it, standing apart enough to take the
gift and, in time, to give in return.

But a fullness with no edge leaves no room for anything else. A separate thing
cannot exist inside a light that fills all space. So the first act was not an
outpouring. It was a \textbf{drawing back}. The source contracted itself, pulling
its light away from a single point and leaving, for the first time, an \textbf{empty
space}. Into that cleared room a separate world could begin. Even in the emptiness
a faint trace of the light remained behind, like a scent in a bottle after it is
poured out.

Into the empty space the source sent a single thin \textbf{ray of light}. The ray
was measured and controlled, unlike the boundless fullness before it. Everything to
come would be built along this one line.

The ray first took the form of a \textbf{template}, the pattern of a whole person
laid over the whole of what would be made. From the openings of this great figure,
its eyes and ears and mouth, the light streamed out, and that streaming became the
worlds.

The light organized into \textbf{ten channels}, ten attributes through which the
hidden source could act and be known. Kindness and limit, wisdom and form, the
will to rule and the readiness to receive. These ten repeated at \textbf{four
levels}, four descending worlds, each one a coarser and more hidden dressing of the
light above it. Reality came down like a chain, link from link, the light growing
dimmer and more clothed at every step.

Then came the break. The early \textbf{vessels}, the containers meant to hold each
of the ten lights, were strong enough at the top but too weak below. They stood
alone, each grasping the full light for itself, with no give-and-take between them.
Strong light in weak, separate containers cannot hold. The lower vessels
\textbf{shattered}. This was the disaster at the root of the world, the source of
all disorder, exile, and death.

When the vessels broke, the shards fell. Clinging to each falling shard was a
\textbf{spark of holy light}, trapped now inside the broken \textbf{husks}.
These husks are the shells that hide the light, the dark side of things. So evil is
not a thing that was made on purpose. It is broken structure with holy light caught
in the wrong place. The light is still holy. It is only imprisoned.

The world we live in is this broken world. It is full of exile and scattered light.

But the work is not finished. The source planted in the broken world a creature
made in the image of the great template, the \textbf{human soul}. By right action
and true intention, the soul can find the trapped sparks, lift them out of the
husks, and send them home. Each good deed gathers a spark. Each gathered spark mends
a piece of the broken vessels.

The story has a clear skeleton. A boundless source. A drawing-back that makes room.
A single ray of first distinction. A descent through worlds into vessels. A break
where the vessels fail. Holy light trapped among the fragments. A self-aware
creature whose labor gathers what was scattered and returns the world to wholeness.

The going-out is creation. The coming-back is history. The one who carries the
world home is the human being.
